\section{Agentic AI 系统抽象与复杂度来源}

课程将 Agentic 系统抽象为四层：感知（Perception）、推理规划（Reasoning/Planning）、工具调用（Tool Use）、执行与反馈（Action/Feedback）。相较于单轮对话模型，Agent 的关键特征是 ``闭环''：动作会改变环境，环境反馈再反作用于下一步决策。

\begin{importantbox}{复杂度不是线性增加，而是组合爆炸}
每增加一个维度（输入模态、动作空间、工具集、记忆类型、自主等级），攻击面不是加法增长，而是组合增长。特别是当工具集从固定白名单扩展到运行时发现时，系统的可攻击状态空间会急速放大。
\end{importantbox}

\begin{table}[H]
\centering
\begin{tabular}{p{2.2cm}p{4.2cm}p{4.5cm}}
\toprule
\textbf{维度} & \textbf{可选范围} & \textbf{安全含义} \\
\midrule
输入空间 & 文本 / 图像 / 文件 / URL / API 返回 & 输入污染、跨模态注入、外部内容投毒 \\
动作空间 & 回答 / 调工具 / 执行命令 / 发请求 & 越权执行、侧向移动、真实系统破坏 \\
工具集合 & 无工具 / 固定工具 / 动态发现工具 & 攻击路径不透明、能力边界漂移 \\
记忆机制 & 无记忆 / 短期记忆 / 持久记忆 & 恶意状态持久化、长期污染 \\
自主等级 & 人工审批 / 半自动 / 全自动 & 风险收敛速度与损害放大倍数不同 \\
\bottomrule
\end{tabular}
\caption{Agent 系统设计维度与安全风险耦合}
\end{table}

\begin{knowledgebox}{Hybrid Agent System 的挑战}
很多现实系统是 ``Neural + Symbolic + Rule'' 混合体：LLM 做规划，规则引擎做审批，工具适配器做执行。攻击者常利用模块边界不一致（policy gap）完成跨层绕过。
\end{knowledgebox}

\subsection{本章小结}
Agent 系统的安全复杂度来自多维组合而非单点漏洞。设计阶段必须把输入、动作、工具、记忆、自主度同时纳入威胁建模。

